\documentclass[10pt,a4paper]{amsart}
\usepackage[margin=19mm]{geometry}
\usepackage{amsmath,amssymb,mathtools}
\usepackage[scr=rsfs,cal=euler]{mathalfa}
\usepackage{xfrac}
\usepackage[unicode,hidelinks,bookmarks=false]{hyperref}
\newcommand{\ie}{i.e.\ }
\newcommand{\cf}{cf.\ }
\newcommand{\resp}{respectively\ }
\newcommand{\Subsection}{\subsection}
\providecommand{\nobreakdash}{\nobreak\hbox{-}}
\setlength{\emergencystretch}{2em}
\multlinegap0pt
\usepackage{amssymb}
\usepackage{dsfont}
\usepackage{tikz}
\newtheorem{problem}{Problem}
\newtheorem{proposition}{Proposition}
\newcommand{\Menge}[2]{\bigg\{{#1}~\bigg|~{#2}\bigg\}} 
\newcommand{\R}{{\mathbb R}}
\newcommand{\minimize}[2]{\ensuremath{\underset{\substack{{#1}}}%
{\mathrm{minimize}}\;\;#2 }}
\newcommand{\reli}{\ensuremath{\text{\rm ri}\,}}
\title{Bregman proximal gradient method for linear optimization under entropic constraints}
\author{Luis M. Brice\~no-Arias, Ma\"el Le Treust}
\date{Source: arXiv:2506.10849v4 (CC BY 4.0)}
\begin{document}
\maketitle

\noindent\emph{Source and licence.} Bregman proximal gradient method for linear optimization under entropic constraints. Version arXiv:2506.10849v4 (CC BY 4.0), distributed by arXiv under the Creative Commons Attribution 4.0 International licence, http://creativecommons.org/licenses/by/4.0/, as recorded in the arXiv licence metadata for this exact version and shown on the abstract page https://arxiv.org/abs/2506.10849v4. Full licence notice: \texttt{licenses/arxiv-2506.10849-CC-BY-4.0.txt}.\par\smallskip

\noindent\textit{Editorial scope.} Counted unit: the article's proposition p:achiev (source0.tex lines 1493--1503). The article's complete original proof of it is delivered with it (source0.tex lines 1505--1527, one \texttt{proof} environment of 23 lines). No mathematics is written, repaired or re-derived: the statement and the proof are the article's own lines, in the article's own notation, and no retained line is substituted or re-flowed.  Retained uncounted as the source-local context the proof needs: lines 446--468; lines 1451--1456.  Nothing was omitted from any retained span, so the package carries no omission record.  No retained line contains a citation, so the wrapper carries no bibliography and no bibliography entry.  No retained line contains a figure environment, an image-inclusion command or drawing code, so no image is delivered and the package declares no asset dependency.  Editorial formatting only, in addition: an AMS \texttt{amsart} preamble that restates the article's own declarations and macros used by the retained lines, namely the environments \texttt{problem}, \texttt{proposition} on separate counters, together with the standard packages those lines need.  The retained environments print in this document as Problem 1, Proposition 1 in order of appearance, whereas the article prints the same items with the numbering of its own full document.  The complete native source of the article as distributed in the arXiv e-print for this exact version is bundled unchanged as \texttt{sources/179757/source0.tex}.
\begin{problem}\label{prob:main}
Let $A$, $B$, and $S$ be nonempty finite sets such that $|A|\ge 
2$, let $(c_{ab}^s)_{a\in A,\, b\in B}^{s\in S}\in \R^{(A\times 
B)^S}$ a vector of cost values, and let $p=(p_s)_{s\in 
S}\in{\reli(\Delta(S))}$ a probability distribution. 
The problem is to
\begin{align}
\label{e:PrimalProblem}
\minimize{q\in C}{\sum_{s\in 
S}p_s\sum_{a\in A}\sum_{b\in B}c_{ab}^sq_{ab}^s},
\end{align}
where 
\begin{equation}
\label{e:defC}
C=\Menge{q\in(\Delta(A\times 
B))^S}{\sum_{s\in 
S}p_s\sum_{a\in A}\sum_{b\in B}q_{ab}^s\ln 
\Bigg(\frac{q_{ab}^s}{\sum\limits_{a'\in A}\sum\limits_{s'\in 
S}p_{s'}q_{a'b}^{s'}}\Bigg)\leq 
0}
\end{equation}
with the convention $0\cdot\ln 0=0$.
\end{problem}

\begin{align}
(\forall s\in S)\quad c_{\textsf{min}}^s &= \min_{(a,b) \in A\times B} 
c_{ab}^s\label{eq:Cmin}\\
c_{\textsf{min}} &= \sum_{s\in S}p_s 
c_{\textsf{min}}^s\label{eq:Cmin2}
\end{align}

\begin{proposition}
\label{p:achiev}
In the context of Problem~\ref{prob:main}, let 
$\overline{q}\in(\Delta(A\times B))^S$. Then 
\begin{equation}
\overline{q}\in\Delta_0\quad \Leftrightarrow\quad 
\min_{q\in(\Delta(A\times B))^S} 
\sum_{a,b,s}p_sc_{ab}^sq_{ab}^s=\sum_{a,b,s}p_sc_{ab}^s
\overline{q}_{ab}^s=c_{\textsf{min}}.
\end{equation}
\end{proposition}

\begin{proof}
Indeed, if $\overline{q}\in\Delta_0$
we have 
\begin{align}
{c_{\textsf{min}}\le }\min_{q\in(\Delta(A\times B))^S} 
\sum_{a,b,s}p_sc_{ab}^sq_{ab}^s 
= \sum_{s\in S}p_s c_{\textsf{min}}^s\!\!\!\!\sum_{(a,b)\in I^s}  
\overline{q}_{ab}^s = \sum_{s\in S}p_s c_{\textsf{min}}^s = 
c_{\textsf{min}}.\label{e:PrimalNoIC}
\end{align}
Conversely, if $\overline{q}\notin\Delta_0$, there exists 
$s\in S$ and $(a,b)\notin I^s$ such that $\overline{q}_{ab}^s>0$.
Since $(a,b)\notin I^s$, we have $c_{ab}^s>c_{\mathsf{min}}^s$ and, 
therefore,
it follows from \eqref{eq:Cmin} that
\begin{equation}
\sum_{a,b,s}p_sc_{ab}^s\overline{q}_{ab}^s>\sum_{s\in S}p_s 
c_{\textsf{min}}^s\sum_{a\in A}\sum_{b\in B}  
\overline{q}_{ab}^s = \sum_{s\in S}p_s c_{\textsf{min}}^s = 
c_{\textsf{min}}
\end{equation}
and the proof is complete.\qed
\end{proof}

\end{document}
